To rigorously investigate the sign of the defect function $\Xi(\theta, \psi)$ over the fundamental sector, we first simplify the sums for $C$ and $D$ by explicitly evaluating the sums over $\epsilon_3 \in \{\pm 1\}$. This leads to a profound algebraic relationship between $C, D$, and $D'$ which ultimately allows us to evaluate the sign analytically.

### 1. Reduction to 4-Term Sums
Let standard spherical coordinates be $x = \sin\theta\cos\psi$, $y = \sin\theta\sin\psi$, $z = \cos\theta$. We define the reduced variables:
$$ u = \frac{x}{\sqrt{3}}, \quad v = \frac{y}{\sqrt{3}}, \quad c = \frac{z}{\sqrt{3}} $$
For a given pair $(\epsilon_1, \epsilon_2) \in \{\pm 1\}^2$, let $b_{\epsilon} = 1 - \epsilon_1 u - \epsilon_2 v$ and $a_{\epsilon} = \epsilon_2 x - \epsilon_1 y$. 
Pairing the terms corresponding to $\epsilon_3 = 1$ and $\epsilon_3 = -1$ in the definitions of $C$ and $D$, we find:
$$ \frac{1}{b_{\epsilon} - c} + \frac{1}{b_{\epsilon} + c} = \frac{2b_{\epsilon}}{b_{\epsilon}^2 - c^2}, \quad \text{and} \quad \frac{1}{b_{\epsilon} - c} - \frac{1}{b_{\epsilon} + c} = \frac{2c}{b_{\epsilon}^2 - c^2} $$
Substituting these pairs into the original 8-term sums yields exactly 4 terms. Let $d_{\epsilon} = b_{\epsilon}^2 - c^2$ and $M_{\epsilon} = \frac{a_{\epsilon}}{d_{\epsilon}}$. With $K_0 = \frac{1}{2\pi\sqrt{3}}$, we obtain:
$$ C(\theta,\psi) = K_0 \sum_{\epsilon \in \{\pm 1\}^2} M_{\epsilon} b_{\epsilon} $$
$$ D(\theta,\psi) = K_0 c \sum_{\epsilon \in \{\pm 1\}^2} M_{\epsilon} \epsilon_1 $$

By definition, $D' = D(\theta, \psi+\pi/2)$ is obtained by rotating $(x,y) \mapsto (-y, x)$. Substituting this into $D$, re-indexing the sum by $(\eta_1, \eta_2) = (\epsilon_2, -\epsilon_1)$, and utilizing $a_{\eta} = -a_{\epsilon}$, we find:
$$ D'(\theta,\psi) = K_0 c \sum_{\epsilon \in \{\pm 1\}^2} M_{\epsilon} (-\epsilon_2) $$

### 2. The Algebraic Identity Relating $C$, $D$, and $D'$
Let us enumerate the four index pairs for $\epsilon = (\epsilon_1, \epsilon_2)$ as $1=(+,+)$, $2=(+,-)$, $3=(-,+)$, and $4=(-,-)$. The terms $b_i$ expand to:
$$ b_1 = 1 - u - v, \quad b_2 = 1 - u + v, \quad b_3 = 1 + u - v, \quad b_4 = 1 + u + v $$
This lets us write $b_{\epsilon} = 1 - \epsilon_1 u - \epsilon_2 v$. Substituting this into the sum for $C$:
$$ C = K_0 \sum_{i=1}^4 M_i (1 - \epsilon_1 u - \epsilon_2 v) = K_0 \left( \sum_{i=1}^4 M_i - u \sum_{i=1}^4 M_i \epsilon_1 - v \sum_{i=1}^4 M_i \epsilon_2 \right) $$
Notice that the second and third sums precisely match our expressions for $D$ and $D'$. Denoting $S = \sum_{i=1}^4 M_i$, we uncover the exact structural identity:
$$ C = K_0 S - \frac{x}{z} D + \frac{y}{z} D' $$

### 3. Explicit Computation of the Residual Sum $S$
The term $S$ reveals the residual behavior of $C$ uncaptured by $D$ and $D'$. Let us evaluate $S = M_1 + M_4 + M_2 + M_3$. We observe that $a_4 = -a_1$ and $a_3 = -a_2$. Hence:
$$ M_1 + M_4 = a_1 \left( \frac{1}{d_1} - \frac{1}{d_4} \right) = a_1 \frac{d_4 - d_1}{d_1 d_4} $$
Since $d_4 - d_1 = b_4^2 - b_1^2 = (b_4-b_1)(b_4+b_1) = 2(u+v) \cdot 2 = 4(u+v)$, and similarly $d_3 - d_2 = 4(u-v)$, we acquire:
$$ S = \frac{4 a_1(u+v)}{d_1 d_4} + \frac{4 a_2(u-v)}{d_2 d_3} $$
Noting $a_1(u+v) = (x-y)\frac{x+y}{\sqrt{3}} = \frac{x^2-y^2}{\sqrt{3}}$ and $a_2(u-v) = -(x+y)\frac{x-y}{\sqrt{3}} = -\frac{x^2-y^2}{\sqrt{3}}$, the numerators are precise opposites:
$$ S = \frac{4(x^2 - y^2)}{\sqrt{3}} \left( \frac{1}{d_1 d_4} - \frac{1}{d_2 d_3} \right) = \frac{4(x^2 - y^2)}{\sqrt{3}} \frac{d_2 d_3 - d_1 d_4}{d_1 d_2 d_3 d_4} $$
A direct algebraic expansion of $d_1 d_4 - d_2 d_3$ leveraging $u^2+v^2+c^2 = \frac{1}{3}$ completely factors it to $-\frac{16xy(1+z^2)}{9}$. This provides a manifestly exact, non-singular form for $S$:
$$ S = -\frac{64 x y (x^2 - y^2) (1+z^2)}{9\sqrt{3} \, d_1 d_2 d_3 d_4} $$

### 4. Rigorous Evaluation of $\Xi$ at the Equator ($z=0$)
To deduce the strict global sign of $\Xi$, we test the expression continuously on the relative interior boundary corresponding to the spherical equator $z=0$ (i.e., $\theta = \pi/2$), which lies well within the fundamental limit of valid topologies.

When $z \to 0$, we clearly have $c \to 0$. From the exact finite sums derived in Step 1, $D$ and $D'$ possess a proportional factor of $c$ in the numerator while their denominators $d_i = b_i^2$ stay bounded away from zero (as $x, y > 0$). Thus, identically:
$$ D(x, y, 0) = 0 \quad \text{and} \quad D'(x, y, 0) = 0 $$
By our algebraic identity, $C(x, y, 0) = K_0 S$. Substituting $z=0$ into our closed-form result for $S$:
$$ C(x, y, 0) = -\frac{32 x y (x^2 - y^2)}{27 \pi \, b_1^2 b_2^2 b_3^2 b_4^2} $$
Within the fundamental sector $(0, \pi) \times (0, \pi/2)$, if we evaluate along the equator, substituting $x = \cos \psi$ and $y = \sin \psi$, we have $x^2 - y^2 = \cos(2\psi)$ and $2xy = \sin(2\psi)$. Consequently, for all $\psi \in (0, \pi/4) \cup (\pi/4, \pi/2)$, we have $C(x, y, 0) \neq 0$. 

Therefore, exactly at the equator subset inside our fundamental domain, we evaluate the defect function:
$$ \Xi(x, y, 0) = \frac{1}{2}(0)^2 + \frac{1}{2}(0)^2 - \left( -\frac{32 x y (x^2 - y^2)}{27 \pi \prod b_i^2} \right)^2 = -C^2 < 0 $$
**Conclusion:** We have rigorously proved that the defect $\Xi(\theta, \psi)$ evaluates to a strictly negative value almost everywhere along the equator line $z=0$ (i.e., $\theta = \pi/2$). This establishes a hard mathematical counterexample demonstrating that $\Xi \ge 0$ does **not** globally hold everywhere strictly within the specified fundamental sector $(0, \pi) \times (0, \pi/2)$, fundamentally necessitating either an adjusted domain or a refactored invariant inequality.